\documentclass[9pt,table]{beamer} % options fleqn

\usepackage{/home/paul/Workspace/build/macros}


\usetheme[titleformat frame = smallcaps]{metropolis}
%\usepackage{FiraSans}
\usepackage{FiraMono}

\metroset{sectionpage=none}

%\usepackage[T1]{fontenc}


\setbeamersize{text margin left=5mm,text margin right=5mm} 
%\setbeamersize{text margin left=9mm,text margin right=3mm} 

\setbeamertemplate{section in toc}[sections numbered]

% algorithms
\usepackage[ruled,linesnumbered]{algorithm2e}
\SetKwInput{KwInput}{Input}
\SetKwInput{KwOutput}{Output}

% typography
\usepackage{soul}

\usepackage{caption}

\usepackage{varwidth}

\usepackage{dsfont}

% figures
\usepackage{svg}
%\svgpath{{figs/}}

% drawings
\usepackage{tikz}
\usetikzlibrary{shapes,shapes.misc,shapes.geometric}
\usetikzlibrary{positioning,fit,backgrounds}
\usetikzlibrary{trees,matrix}
\usetikzlibrary{decorations.pathreplacing,calligraphy}
%
\tikzset{cross/.style={cross out, draw=red, minimum size=3pt, line width=0.25mm, inner sep=0pt, outer sep=0pt}, cross/.default={1pt}}
%
\usepackage[skins,theorems]{tcolorbox}
\tcbset{highlight math style={colback=mLightBrown!20,colframe=black!2,sharp corners},bottom=2pt,top=2pt}
\tcbset{left skip={-.01\linewidth}, colback=mLightBrown!20,colframe=black!2,sharp corners,left=0pt,right=0pt}
%
\newcommand{\point}[1]{\textbf{\textcolor{mDarkTeal}{#1}}\quad}
%
%\definecolor{cobalt}{HTML}{0047AB}
\definecolor{aqua}{HTML}{00FFFF}
\newcommand{\kw}[1]{\textbf{\alert{#1}}}
\newcommand{\invsible}[1]{\textcolor{black!2}{#1}}
%
\definecolor{myteal}{rgb}{0.21, 0.46, 0.53}

\usepackage[position=bottom]{subfig}

\usepackage{tabularx}
\usepackage{booktabs}% http://ctan.org/pkg/booktabs
\newcommand{\tabitem}{\llap{\tiny$\blacksquare$}~~}

% handling lists
\usepackage{enumitem}

\usepackage{multimedia}

% tcolorbox
\usepackage[skins,theorems]{tcolorbox}
\tcbset{highlight math style={colback=mLightBrown!20,colframe=black!2,sharp corners},bottom=2pt,top=2pt}
\tcbset{left skip={-.01\linewidth}, colback=mLightBrown!20,colframe=black!2,sharp corners,left=0pt,right=0pt}


% references
\usepackage[backend=biber,style=authoryear]{biblatex}
%\usepackage[backend=biber,style=alphabetic]{biblatex}
\bibliography{/home/paul/Workspace/build/biblio}
\usepackage{bibentry}
\newcommand{\mycite}[2]{\footnotetext[#1]{\emph{\tiny{\citeauthor{#2},\,\citeyear{#2},}}\,\tiny{\citefield{#2}{title}}}}

%\renewcommand{\cite}[1]{\textcolor{myteal}{[\citeauthor{#1},\,\citeyear{#1}]}}%
%\newcommand{\mycite}[1]{\textcolor{myteal}{\cite{#1}}}%
\AtBeginEnvironment{frame}{\setcounter{footnote}{0}}

\usepackage{hyperref}


\graphicspath{{figs/}{examples/}{/home/paul/Workspace/data/standard_test_images}}




% add a sommaire
%\makeatletter
%\setbeamertemplate{headline}{%
%  \begin{beamercolorbox}[colsep=1.5pt]{upper separation line head}
%  \end{beamercolorbox}
%  \begin{beamercolorbox}{section in head/foot}
%    \vskip2pt\insertnavigation{\paperwidth}\vskip2pt
%  \end{beamercolorbox}%
%  \begin{beamercolorbox}[colsep=1.5pt]{lower separation line head}
%  \end{beamercolorbox}
%}
%\makeatother

\setbeamertemplate{theorem begin}
{%
  \inserttheoremheadfont% \bfseries
  \inserttheoremname \inserttheoremnumber
  \ifx\inserttheoremaddition\@empty\else\ (\inserttheoremaddition)\fi%
  \inserttheorempunctuation
  \normalfont
}
\setbeamertemplate{theorem end}{%
  % empty
}
\makeatother

\setbeamercolor{section in head/foot}{fg=normal text.bg, bg=structure.fg}



\title{Optimisation Numérique et Sciences des Données}
\subtitle{Sciences des données -- Introduction aux réseaux de neurones}
\author{Paul Catala (CRAN, UL) \hfill Master Ingénierie des Systèmes Complexes M2 ISC 925}
%\date{Postdoctorant à l'université d'Osnabrück, Allemagne}
\date{paul.catala(at)univ-lorraine.fr \hfill Semestre 9 -- 2024 \\}

%\usepackage{uoslogo}
%\titlegraphic{\uoslogored{4cm}}



\begin{document}
\maketitle



\section{Introduction}

\begin{frame}
\frametitle{Programme du cours}
\tableofcontents[currentsection]
\end{frame}

\begin{frame}
\frametitle{Contexte: apprentissage statistique}

\begin{itemize}[label=\tiny$\blacksquare$]
	\item Données \alert{labellisées}
	\eq{
		(\xx_i,y_i) \in \RR^n \times \Yy
	}
	
	\item \alert{Apprentissage supervisé}: trouver une fonction \alert{paramétrée} $f_\theta : \RR^n \to \Yy$ t.q.
	\eq{
		f_\theta(\xx_i) = y_i
	}
	où $\theta$ est un vecteur des \alert{paramètres} (qui caractérisent le classifieur).
	
	\pause
	\item Pour des tâches de classification binaire ($\Yy = \{0,1\}$), nous avons vu les \alert{machines à vecteurs de support}: étant donnée une transformation $\phi:\RR^n \to \RR^m$,
	\eq{
		f_\theta(\xx) = \sgn(\abold^\top \phi(\xx) + b ) \qquad (\theta = (\abold,b) \in \RR^{m+1})
	}
	où le $\theta$ est choisi de manière à \alert{maximiser la marge} par rapport aux données (transformées) $\phi(\xx_i)$
		
\end{itemize}
\end{frame}


\begin{frame}
\frametitle{Contexte: apprentissage statistique}

\begin{itemize}[label=\tiny$\blacksquare$]
	\item L'apprentissage statistique ne se réduit pas à de la classification binaire supervisée
	
	\item \textbf{Classication supervisée, semi-supervisée, non supervisée}: méthode des K plus proches voisins ($K$-means), forêt d'arbres décisionnels (random forest), SVM multi-classes, \alert{réseaux de neurones}
	
	\item \textbf{Régression}
\end{itemize}
\end{frame}




\begin{frame}
\frametitle{Réseau de neurones}

\newlength{\spacing}
\setlength{\spacing}{1.5cm}

\vspace{-1cm}

\centering

\begin{tikzpicture}[every node/.style={circle,minimum width = .5cm}]

% first layer
\node[draw] (n10) at (0,0) {$x_1$};
\node[draw,below=5pt of n10] (n11) {$x_2$};
\node[draw,below=5pt of n11] (n12) {$x_3$};
\node[draw,below=5pt of n12] (n13) {$x_4$};


% second layer
\node[draw,above right = 0.77cm and \spacing of n10] (n20) {};
\node[draw,below=5pt of n20] (n21) {};
\node[draw,below=5pt of n21] (n22) {};
\node[draw,below=5pt of n22] (n23) {};
\node[draw,below=5pt of n23] (n24) {};
\node[draw,below=5pt of n24] (n25) {};
\node[draw,below=5pt of n25] (n26) {};
\node[draw,below=5pt of n26] (n27) {};

% connections
\draw[] (n10) -- (n20);
\draw[] (n10) -- (n21);
\draw[] (n10) -- (n22);
\draw[] (n10) -- (n23);
\draw[] (n10) -- (n24);
\draw[] (n10) -- (n25);
\draw[] (n10) -- (n26);
\draw[] (n10) -- (n27);
%
\draw[] (n11) -- (n20);
\draw[] (n11) -- (n21);
\draw[] (n11) -- (n22);
\draw[] (n11) -- (n23);
\draw[] (n11) -- (n24);
\draw[] (n11) -- (n25);
\draw[] (n11) -- (n26);
\draw[] (n11) -- (n27);
%
\draw[] (n12) -- (n20);
\draw[] (n12) -- (n21);
\draw[] (n12) -- (n22);
\draw[] (n12) -- (n23);
\draw[] (n12) -- (n24);
\draw[] (n12) -- (n25);
\draw[] (n12) -- (n26);
\draw[] (n12) -- (n27);
%
\draw[] (n13) -- (n20);
\draw[] (n13) -- (n21);
\draw[] (n13) -- (n22);
\draw[] (n13) -- (n23);
\draw[] (n13) -- (n24);
\draw[] (n13) -- (n25);
\draw[] (n13) -- (n26);
\draw[] (n13) -- (n27);


% third layer
\node[draw,below right = 0.3cm and \spacing of n20] (n30) {};
\node[draw,below=5pt of n30] (n31) {};
\node[draw,below=5pt of n31] (n32) {};
\node[draw,below=5pt of n32] (n33) {};
\node[draw,below=5pt of n33] (n34) {};
\node[draw,below=5pt of n34] (n35) {};

% connections
\draw[] (n20) -- (n30);
\draw[] (n20) -- (n31);
\draw[] (n20) -- (n32);
\draw[] (n20) -- (n33);
\draw[] (n20) -- (n34);
\draw[] (n20) -- (n35);
%
\draw[] (n21) -- (n30);
\draw[] (n21) -- (n31);
\draw[] (n21) -- (n32);
\draw[] (n21) -- (n33);
\draw[] (n21) -- (n34);
\draw[] (n21) -- (n35);
%
\draw[] (n22) -- (n30);
\draw[] (n22) -- (n31);
\draw[] (n22) -- (n32);
\draw[] (n22) -- (n33);
\draw[] (n22) -- (n34);
\draw[] (n22) -- (n35);
%
\draw[] (n23) -- (n30);
\draw[] (n23) -- (n31);
\draw[] (n23) -- (n32);
\draw[] (n23) -- (n33);
\draw[] (n23) -- (n34);
\draw[] (n23) -- (n35);
%
\draw[] (n24) -- (n30);
\draw[] (n24) -- (n31);
\draw[] (n24) -- (n32);
\draw[] (n24) -- (n33);
\draw[] (n24) -- (n34);
\draw[] (n24) -- (n35);
%
\draw[] (n25) -- (n30);
\draw[] (n25) -- (n31);
\draw[] (n25) -- (n32);
\draw[] (n25) -- (n33);
\draw[] (n25) -- (n34);
\draw[] (n25) -- (n35);
%
\draw[] (n26) -- (n30);
\draw[] (n26) -- (n31);
\draw[] (n26) -- (n32);
\draw[] (n26) -- (n33);
\draw[] (n26) -- (n34);
\draw[] (n26) -- (n35);
%

\draw[] (n27) -- (n30);
\draw[] (n27) -- (n31);
\draw[] (n27) -- (n32);
\draw[] (n27) -- (n33);
\draw[] (n27) -- (n34);
\draw[] (n27) -- (n35);


% fourth layer
\node[draw,above right = 0.3cm and \spacing of n30] (n40) {};
\node[draw,below=5pt of n40] (n41) {};
\node[draw,below=5pt of n41] (n42) {};
\node[draw,below=5pt of n42] (n43) {};
\node[draw,below=5pt of n43] (n44) {};
\node[draw,below=5pt of n44] (n45) {};
\node[draw,below=5pt of n45] (n46) {};
\node[draw,below=5pt of n46] (n47) {};


% connections
\draw[] (n30) -- (n40);
\draw[] (n30) -- (n41);
\draw[] (n30) -- (n42);
\draw[] (n30) -- (n43);
\draw[] (n30) -- (n44);
\draw[] (n30) -- (n45);
\draw[] (n30) -- (n46);
\draw[] (n30) -- (n47);
%
\draw[] (n31) -- (n40);
\draw[] (n31) -- (n41);
\draw[] (n31) -- (n42);
\draw[] (n31) -- (n43);
\draw[] (n31) -- (n44);
\draw[] (n31) -- (n45);
\draw[] (n31) -- (n47);
%
\draw[] (n32) -- (n40);
\draw[] (n32) -- (n41);
\draw[] (n32) -- (n42);
\draw[] (n32) -- (n43);
\draw[] (n32) -- (n44);
\draw[] (n32) -- (n45);
\draw[] (n32) -- (n46);
\draw[] (n32) -- (n47);
%
\draw[] (n33) -- (n40);
\draw[] (n33) -- (n41);
\draw[] (n33) -- (n42);
\draw[] (n33) -- (n43);
\draw[] (n33) -- (n44);
\draw[] (n33) -- (n45);
\draw[] (n33) -- (n46);
\draw[] (n33) -- (n47);
%
\draw[] (n34) -- (n40);
\draw[] (n34) -- (n41);
\draw[] (n34) -- (n42);
\draw[] (n34) -- (n43);
\draw[] (n34) -- (n44);
\draw[] (n34) -- (n45);
\draw[] (n34) -- (n46);
\draw[] (n34) -- (n47);
%
\draw[] (n35) -- (n40);
\draw[] (n35) -- (n41);
\draw[] (n35) -- (n42);
\draw[] (n35) -- (n43);
\draw[] (n35) -- (n44);
\draw[] (n35) -- (n45);
\draw[] (n35) -- (n46);
\draw[] (n35) -- (n47);
%

% output layer
\node[draw,below right = 2cm and \spacing of n40] (n50) {$y$};
%\node[below=5pt of n50] (n51) {$y_2$};
%\node[below=5pt of n51] (n52) {$y_3$};


% connections
\draw[] (n40) -- (n50);
%\draw[] (n40) -- (n51);
%\draw[] (n40) -- (n52);
%
\draw[] (n41) -- (n50);
%\draw[] (n41) -- (n51);
%\draw[] (n41) -- (n52);
%
\draw[] (n42) -- (n50);
%\draw[] (n42) -- (n51);
%\draw[] (n42) -- (n52);
%
\draw[] (n43) -- (n50);
%\draw[] (n43) -- (n51);
%\draw[] (n43) -- (n52);
%
\draw[] (n44) -- (n50);
%\draw[] (n44) -- (n51);
%\draw[] (n44) -- (n52);
%
\draw[] (n45) -- (n50);
%\draw[] (n45) -- (n51);
%\draw[] (n45) -- (n52);
%
\draw[] (n46) -- (n50);
%\draw[] (n46) -- (n51);
%\draw[] (n46) -- (n52);
%
\draw[] (n47) -- (n50);
%\draw[] (n47) -- (n51);
%\draw[] (n47) -- (n52);


\begin{scope}[on background layer]
\node [rectangle, fill=mLightBrown!20, minimum height = 5.6cm, fit={(n10) (n11) (n12) (n13) }, label={[label distance = -23pt]\color{mLightBrown}Couche d'entrée}] (layer-1) {};

\node [rectangle, fill=teal!20, minimum height = 5.6cm, fit={(n20) (n21) (n22) (n23) (n24) (n25) (n26) (n27)}] (layer-2) {};

\node [rectangle, fill=teal!20, minimum height = 5.6cm, fit={(n30) (n31) (n32) (n33) (n34) (n35) }, label={[label distance = -25pt]\color{teal}Couches latentes}] (layer-3) {};

\node [rectangle, fill=teal!20, minimum height = 5.6cm, fit={(n40) (n41) (n42) (n43) (n44) (n45) (n46) (n47)}] (layer-4) {};

\node [rectangle, fill=purple!20, minimum height = 5.6cm, fit={(n50)}, label={[label distance = -25pt]\color{purple}Couche de sortie}] (layer-5) {};
\end{scope}

\draw[-latex,thick] ([yshift=-.3cm]layer-1.south west) -- ([yshift=-.3cm]layer-5.south east) node[rectangle,midway,fill=black!2] {$y = f_\theta(x)$};

\end{tikzpicture}

Objectif: minimiser l'erreur $\sum_j \ell(f_\theta(x^{(j)}),y^{(j)})$ sur des données $(x^{(j)},y^{(j)})$

Variables: $\sim 10^6$ à $10^9$ paramètres $\theta$

\end{frame}






%\bibliographystyle{apalike}
%\printbibliography
%\bibliography{/export/home/pcatala/Workspace/build/biblio}

\end{document}
